\section{Why Stop-Loss Calibration, Not Plug-In Mean Calibration}
\label{app:calibrator-target}
\label{app:stoploss-vs-plugin}

CVaR is nonlinear in the outcome distribution. This appendix explains the
design choice that follows: calibrate shortfalls directly, rather than taking
the tail of calibrated mean predictions.

\subsection{Two Candidate Targets}

The stop-loss route fits one calibrator per threshold:
\begin{equation}
\widehat V^{\mathrm{SL}}_\alpha
=\max_t\left[
t-\frac{1}{\alpha}\E_\ptarget[\hat h_t(\Sscore,\X)]
\right],
\qquad
\hat h_t(\Sscore,\X)\approx
\E[\shortfall{t}{\Y}\mid\Sscore,\X].
\label{eq:app-stoploss}
\end{equation}
The plug-in route fits one mean calibrator and takes the tail of its
predictions:
\begin{equation}
\widehat V^{\mathrm{PI}}_\alpha
=\max_t\left[
t-\frac{1}{\alpha}
\E_\ptarget[\shortfall{t}{\hat m(\Sscore,\X)}]
\right],
\qquad
\hat m(\Sscore,\X)\approx\E[\Y\mid\Sscore,\X].
\label{eq:app-plugin}
\end{equation}
The two look similar, but their population targets are different.

\begin{proposition}[Plug-in mean calibration targets the wrong tail]
\label{prop:plugin-target}
At the population limit, the stop-loss route targets $\CVaR_\alpha(\Y)$ under
shortfall transport. The plug-in route targets
$\CVaR_\alpha(\E[\Y\mid\Sscore,\X])$. These are equal only under degenerate
residual uncertainty, or at thresholds where all conditional outcome mass lies
on one side of the threshold.
\end{proposition}

\begin{twocolproof}
\twocolrow{
\[
\hat h_t(\Sscore,\X)\to
\E[\shortfall{t}{\Y}\mid\Sscore,\X].
\]
}{
This is the population target of the stop-loss calibrator.
}
\twocolrow{
\[
\E[\hat h_t(\Sscore,\X)]
\to \E[\shortfall{t}{\Y}].
\]
}{
Apply iterated expectations.
}
\twocolrow{
\[
\max_t\left[t-\alpha^{-1}\E[\shortfall{t}{\Y}]\right]
=\CVaR_\alpha(\Y).
\]
}{
Use the stop-loss representation in \cref{eq:app-cvar-definition}.
}
\twocolrow{
\[
\hat m(\Sscore,\X)\to\E[\Y\mid\Sscore,\X].
\]
}{
This is the population target of the mean calibrator.
}
\twocolrow{
\[
\max_t\left[t-\alpha^{-1}\E[
\shortfall{t}{\hat m(\Sscore,\X)}]\right]
\to
\CVaR_\alpha(\E[\Y\mid\Sscore,\X]).
\]
}{
The plug-in route computes CVaR on the distribution of conditional means, not
on the distribution of realized oracle outcomes.
}
\end{twocolproof}

\subsection{Jensen Gap}

For fixed $t$, the map $u\mapsto(t-u)_+$ is convex. Therefore
\begin{equation}
\shortfall{t}{\E[\Y\mid\Sscore,\X]}
\le
\E[\shortfall{t}{\Y}\mid\Sscore,\X].
\label{eq:jensen-shortfall}
\end{equation}
After taking expectations, the plug-in objective is at least as large as the
stop-loss objective at every threshold:
\[
t-\alpha^{-1}\E[\shortfall{t}{\E[\Y\mid\Sscore,\X]}]
\ge
t-\alpha^{-1}\E[\shortfall{t}{\Y}].
\]
For a lower-tail score, larger CVaR is less pessimistic. Thus the plug-in
route systematically smooths away residual tail risk whenever the cheap score
does not determine the oracle outcome.

\begin{remark}[Why more data does not fix plug-in CVaR]
The gap in \cref{eq:jensen-shortfall} is a target mismatch, not a finite-sample
error. More calibration data makes $\hat m$ converge more accurately to
$\E[\Y\mid\Sscore,\X]$, which is the wrong object for CVaR unless the
conditional residual variance is zero.
\end{remark}

\subsection{Validation Design}

The controlled DGP used to validate this design choice is intentionally simple:
\[
\Sscore\sim\mathrm{Uniform}(0,1),\qquad
\Y=\mathrm{clip}(0.2+0.6\Sscore+\varepsilon,0,1),
\qquad
\varepsilon\sim N(0,\sigma^2).
\]
At $\sigma=0$, the cheap score determines the oracle outcome and the two
routes agree. As $\sigma$ grows, the residual conditional variance grows while
$\E[\Y\mid\Sscore]$ stays fixed. The expected diagnostic pattern is:
\begin{itemize}[leftmargin=*,nosep]
\item the stop-loss estimator tracks $\CVaR_\alpha(\Y)$;
\item the plug-in estimator tracks $\CVaR_\alpha(\E[\Y\mid\Sscore])$;
\item the plug-in gap grows with $\sigma$ and is larger for smaller $\alpha$.
\end{itemize}
This validation belongs in the appendix because it justifies the estimator's
regression target. It should not be used as a headline empirical result.
